% Part: first-order-logic
% Chapter: syntax-and-semantics
% Section: introduction

\documentclass[../../../include/open-logic-section]{subfiles}

\begin{document}

\olfileid{fol}{syn}{int}

\olsection{مقدمه}

برای بسط نظریه و فرانظریهٔ منطق مرتبهٔ اول، نخست باید نحو و معناشناسی
عبارت‌های آن را تعریف کنیم. عبارت‌های منطق مرتبهٔ اول ترم‌ها و
!!{formula}ها هستند. ترم‌ها از !!{variable}ها، !!{constant}ها و
!!{function}ها ساخته می‌شوند. !!^{formula}ها نیز از !!{predicate}ها
همراه با ترم‌ها ساخته می‌شوند؛ این‌ها کوچک‌ترین !!{formula}ها، یعنی
فرمول‌های «اتمی»، را تشکیل می‌دهند. سپس می‌توانیم با رابط‌های
منطقی و سورها، از !!{formula}های اتمی فرمول‌های پیچیده‌تری بسازیم.
راه‌های گوناگونی برای بیان قواعد ساخت وجود دارد و ما فقط یکی از
امکان‌ها را ارائه می‌کنیم. دستگاه‌های دیگر نمادهای متفاوتی برمی‌گزینند،
مجموعه‌های متفاوتی از رابط‌ها را اولیه می‌گیرند، و پرانتزها را به
شیوه‌ای دیگر به کار می‌برند یا، مانند نمادگذاری موسوم به لهستانی، اصلاً
به کار نمی‌برند. بااین‌حال وجه مشترک همهٔ رویکردها این است که قواعد
ساخت، مجموعهٔ ترم‌ها و !!{formula}ها را \emph{به‌طور استقرایی} تعریف
می‌کنند. اگر این کار درست انجام شود، هر عبارت بر پایهٔ قواعد ساخت
اساساً فقط به یک شیوه ساخته می‌شود. چنین تعریف استقرایی‌ای عبارت‌هایی
با \emph{خوانش یکتا} تولید می‌کند؛ ازاین‌رو می‌توانیم معنای آن‌ها را
نیز با همان روش، یعنی تعریف استقرایی، تعیین کنیم.

تعیین معنای عبارت‌ها موضوع معناشناسی است. مفهوم محوری معناشناسی ارضا
در !!a{structure} است. !!^a{structure} به اجزای سازندهٔ زبان معنا
می‌دهد: !!a{domain} مجموعه‌ای ناتهی از اشیاست. سورها چنان تفسیر
می‌شوند که بر این دامنه تغییر کنند؛ به !!{constant}ها عناصری از دامنه
نسبت داده می‌شود؛ به نمادهای !!{function}ی $n$موضعی تابع‌هایی از
$D^n$ به $D$، که $D$ همان !!{domain} است، نسبت داده می‌شود؛ و به
!!{predicate}ها رابطه‌هایی بر !!{domain} نسبت داده می‌شود. این
!!{domain} همراه با انتساب‌ها به واژگان پایه، !!a{structure} را تشکیل
می‌دهد.
!!^{variable}ها ممکن است در !!{formula}ها ظاهر شوند و برای ارائهٔ
معناشناسی باید !!{element}هایی از !!{domain} را نیز به آن‌ها نسبت
دهیم؛ این همان انتساب متغیر است. در پایان، رابطهٔ ارضا این اجزا را به
هم پیوند می‌دهد. ممکن است !!^a{formula} در !!a{structure}~$\Struct{M}$
نسبت به انتساب متغیر~$s$ ارضا شود، که آن را
$\Sat{M}{!A}[s]$ می‌نویسیم. این رابطه نیز با استقرا بر ساختار~$!A$
تعریف می‌شود؛ برای نمونه، جدول‌های ارزش رابط‌های منطقی، ارضای
$!A \land !B$ را بر حسب ارضا یا عدم ارضای $!A$ و~$!B$ تعریف می‌کنند.
سپس معلوم می‌شود اگر !!{formula}~$!A$ یک !!a{sentence} باشد، یعنی هیچ
متغیر آزادی نداشته باشد، انتساب متغیر بی‌اهمیت است؛ بنابراین می‌توانیم
صرفاً از ارضاشدن یا ارضانشدن !!{sentence}ها در !!{structure}ها سخن
بگوییم.

بر پایهٔ رابطهٔ ارضای $\Sat{M}{!A}$ برای جمله‌ها می‌توانیم مفاهیم
معناشناختی بنیادی اعتبار، پیامد و ارضاپذیری را تعریف کنیم. یک جمله
معتبر است، $\Entails !A$، اگر هر ساختاری آن را ارضا کند. جمله از
مجموعه‌ای از !!{sentence}ها نتیجه می‌شود، $\Gamma \Entails !A$، اگر هر
!!{structure} که همهٔ !!{sentence}های~$\Gamma$ را ارضا می‌کند، $!A$ را
نیز ارضا کند. مجموعه‌ای از جمله‌ها ارضاپذیر است اگر !!{structure}ی
وجود داشته باشد که همهٔ !!{sentence}های آن را هم‌زمان ارضا کند. چون
!!{formula}ها به‌طور استقرایی تعریف شده‌اند و ارضا نیز به نوبهٔ خود با
استقرا بر ساختار !!{formula}ها تعریف می‌شود، می‌توانیم برای اثبات
ویژگی‌های معناشناسی و ربط‌دادن مفاهیم معناشناختی تعریف‌شده از استقرا
استفاده کنیم.

\end{document}
